\chapter{Proofs for Thermodynamic Fisher Geometry}
\label{app:thermodynamic-fisher-draft-proofs}

This appendix proves the statements of \cref{ch:thermodynamic-fisher-draft}
under \cref{hyp:rg-thermo-native-source-box}, with the three families of
\cref{def:rg-thermo-native-families}, and states the lattice-side label
results that accompany the finite-mixture Fisher asymptotics of the main
manuscript. Every volume limit is taken at a fixed word.

\section{Native response and the finite-free boundary transfer}
\label{sec:rg-thermo-native-proof}

The response construction of
\cref{sec:rg-spatial-proof-constrained,sec:rg-spatial-proof-fisher} uses
finite all-touching specifications and marginals of the infinite law. The
finite free law instead omits the plaquettes that leave the square.
\Cref{lem:rg-thermo-native-free-contraction} shows that the estimates of
those sections hold for the free law, and
\cref{lem:rg-thermo-native-local-response,lem:rg-thermo-native-moving-response,lem:rg-thermo-native-row-jets}
prove what is new, the comparison of finite free and infinite quantities.
The free law of a canonical coarse interaction is a fourth family and does
not enter.

\subsection{The free-boundary transfer}
\label{sec:rg-thermo-native-free-response}

Throughout, $\mathfrak D$ is a finite whole-block inventory inside the fine
square of a finite free law, with a prescribed-or-absent pattern of path
parities and a fixed exterior completion, and its constrained law is the free
law conditioned accordingly. For a target block $i$ and an exterior fine edge
$e$, $n(i,e)$ counts the plaquettes meeting both; $\operatorname{cum}$ is the
joint cumulant of bounded observables; $\mathsf T_{i,s}$ is the edge or
plaquette source of the natural derivative direction $i\in\{h,j\}$ at the
site $s$; $Q_R$ is the square of cells within Chebyshev distance $R$ of a
target; $\nu^{\mathfrak D}_{\mathsf g}$ is the finite constrained law; and for
a target cell and a past $S\subset\mathfrak D$ preceding it in a fixed
lexicographic order, $p^{\mathfrak D}_S(\mathsf g,s;y)$ is the four-state row
of the push conditional on $Y_S=y_S$, with $l=\log p^{\mathfrak D}_S$ and
subscripts denoting natural derivatives.

\lemmaheading{Free-boundary transfer}{lem:rg-thermo-native-free-contraction}
The following hold with the constants
\cref{eq:rg-thermo-native-compact-constants}, uniformly in the inventory, the
pattern, the completion, and $\mathsf g\in\widehat{\mathcal K}$.
Changing one incident plaquette moves the field of one coordinate from
$\bar h+j$ to $\bar h-j$ for a baseline field $\bar h$, and the marginal of
that coordinate changes in total variation by
\begin{equation}
 \frac{\sinh(2j)}{\cosh(2\bar h)+\cosh(2j)}\le\tanh j\le\tanh j_+.
 \label{eq:rg-thermo-native-fiber-tilt}
\end{equation}
For a target anchor set $A$, weights
$w_i=e^{-\tau\operatorname{dist}_\infty(\operatorname{anchor}(i),A)}$, and the
expected block Hamming discrepancies $D_i$ of a stationary coupling of two
constrained laws with exteriors $\eta,\eta'$,
\begin{equation}
 S\le\mathrm qS+\mathrm q\sum_{e\in\partial_f\mathfrak D}w_{\operatorname{block}(e)}\delta_e,
 \qquad S=\sum_{i\text{ free}}w_iD_i,
 \quad \delta_e=\mathbf1_{\{\eta_e\ne\eta'_e\}},
 \label{eq:rg-thermo-native-weighted-comparison}
\end{equation}
and for a vector $\mathbf u$ supported on $A$ and a vector $\mathbf w$ outside its radius-$R$
neighborhood,
\begin{equation}
 \operatorname{TV}(\mathcal L(\mathbf u,\mathbf w),\mathcal L(\mathbf u)\otimes\mathcal L(\mathbf w))
 \le\min\{1,\mathrm c_1|A|B^2(R+1)e^{-\tau(R+1)}\}.
 \label{eq:rg-thermo-native-group-decoupling}
\end{equation}
The cumulants satisfy
\begin{equation}
 \begin{gathered}
 \operatorname{cum}(F_0,\ldots,F_r)
 =\sum_{\pi}(-1)^{|\pi|-1}(|\pi|-1)!
                  \prod_{D\in\pi}\E\prod_{i\in D}F_i,\\
 (k_2,\bar k_2)=(2,6),\qquad (k_3,\bar k_3)=(6,26),\qquad
 (k_4,\bar k_4)=(26,150),
 \end{gathered}
 \label{eq:rg-thermo-native-cumulant-partitions}
\end{equation}
where $k_m=\sum_\pi(|\pi|-1)!$ bounds an individual cumulant and
$\bar k_m=2\sum_\pi|\pi|!$ multiplies the total variation of a change of the joint
law. For a bounded parameter-independent observable $f$ with anchor support
$A$, $a_0=|A|\ge1$, and $1\le r\le3$, put
\begin{equation}
 \begin{split}
 L_r(a_0)&=e^{4\tau}[k_{r+1}+\bar k_{r+1}\mathrm c_1(a_0+4r)],\\
 S_r(a_0)&=\frac{(8a_0)^rL_r(a_0)(2r+1)!}
                    {(1-e^{-\tau/r})^{2r+2}},\qquad
 U_r(B,a_0)=S_r(a_0)B^{2r+2},\quad U_0(B,a_0)=1;
 \end{split}
 \label{eq:rg-thermo-native-response-constants}
\end{equation}
then, for ordered sources of the prescribed directions,
\begin{equation}
 \begin{split}
 \sum_{s_1,\ldots,s_r}
 |\operatorname{cum}(f,\mathsf T_{i_1,s_1},\ldots,\mathsf T_{i_r,s_r})|
 &\le U_r(B,a_0)\|f\|_\infty,\\
 \sum_{\max\operatorname{dist}(x_s,A)>n}
 |\operatorname{cum}(f,\mathsf T_{i_1,s_1},\ldots,\mathsf T_{i_r,s_r})|
 &\le U_r(B,a_0)(n+2)^{2r+1}e^{-\tau(n+1)/r}\|f\|_\infty,
 \end{split}
 \label{eq:rg-thermo-native-source-tails}
\end{equation}
and a local $C^2$ observable $f_{\mathsf g}$ has, for $|\alpha|\le2$,
\begin{equation}
 \partial^\alpha\E_{\mathsf g} f_{\mathsf g}=
 \sum_{\beta\le\alpha}\binom{\alpha}{\beta}
 \sum_{s_1,\ldots,s_{|\alpha-\beta|}}
 \operatorname{cum}_{\mathsf g}(\partial^\beta f_{\mathsf g},\mathsf T_{i_1,s_1},\ldots,
                                  \mathsf T_{i_{|\alpha-\beta|},s_{|\alpha-\beta|}}),
 \label{eq:rg-thermo-native-moving-response}
\end{equation}
the source directions being those of $\alpha-\beta$ and an empty source list
meaning expectation. The rows are positive and $C^3$, with
\begin{equation}
 p^{\mathfrak D}_S(\mathsf g,s;y)\ge\frac{(1-\rho_*^B)^2}{4}\ge\frac{d_*^2}{4};
 \label{eq:rg-thermo-native-row-positivity}
\end{equation}
every order-$n$ derivative of a target parity mean is bounded by $M_n(B)$,
where
\begin{equation}
 \begin{split}
 D_n&=g_n+\sum_{k=0}^{n-1}\binom{n}{k}g_kS_{n-k}(9),\qquad g_n=2^nn!,\\
 M_0(B)&=\rho_*^B,\qquad
 M_n(B)=D_nB^{2n+2}\rho_*^{B-n},\qquad 1\le n\le3;
 \end{split}
 \label{eq:rg-thermo-native-parity-jets}
\end{equation}
the logarithmic derivatives of one Bernoulli factor $(1+s m)/2$ have
amplitudes
\begin{equation}
 A_1=\frac{M_1}{d_*},\qquad
 A_2=\frac{M_2}{d_*}+\frac{M_1^2}{d_*^2},\qquad
 A_3=\frac{M_3}{d_*}+\frac{3M_2M_1}{d_*^2}+\frac{2M_1^3}{d_*^3},
 \label{eq:rg-thermo-native-log-amplitudes}
\end{equation}
so that the order-$n$ derivatives of $l$ are bounded by $\widehat A_n=2A_n$;
the row-Fisher observable
$\mathsf I^{\mathfrak D}_{S,ik}=\sum_sp^{\mathfrak D}_Sl_il_k$ has, at fixed $y$,
\begin{equation}
 \begin{split}
 \partial_m\mathsf I^{\mathfrak D}_{S,ik}
 &=\sum_s p^{\mathfrak D}_S(l_ml_il_k+l_{im}l_k+l_il_{km}),\\
 \partial_n\partial_m\mathsf I^{\mathfrak D}_{S,ik}
 &=\sum_s p^{\mathfrak D}_S\bigl[(l_{mn}+l_ml_n)l_il_k
 +l_m(l_{in}l_k+l_il_{kn})+l_n(l_{im}l_k+l_il_{km})\\
 &\quad+l_{imn}l_k+l_{im}l_{kn}+l_{in}l_{km}+l_il_{kmn}\bigr],
 \end{split}
 \label{eq:rg-thermo-native-row-fisher-jets}
\end{equation}
with entries and first two derivatives bounded by
\begin{equation}
 E_B^F=\max\{\widehat A_1^2,\widehat A_1^3+2\widehat A_1\widehat A_2,
          \widehat A_1^4+5\widehat A_1^2\widehat A_2+2\widehat A_1\widehat A_3+2\widehat A_2^2\},
 \qquad H_B^F=\max(1,\widehat A_1,\widehat A_2,\widehat A_3);
 \label{eq:rg-thermo-native-row-fisher-bound}
\end{equation}
if two target parity means have jet differences bounded by
$v=(v_0,v_1,v_2,v_3)$, the one-factor log-jet differences are bounded by
\begin{equation}
 \begin{split}
 \mathcal T_1(v)&=v_1/d_*+M_1v_0/d_*^2,\\
 \mathcal T_2(v)&=v_2/d_*+(M_2v_0+2M_1v_1)/d_*^2+2M_1^2v_0/d_*^3,\\
 \mathcal T_3(v)&=v_3/d_*+(M_3v_0+3M_2v_1+3M_1v_2)/d_*^2\\
 &\quad +(6M_2M_1v_0+6M_1^2v_1)/d_*^3+6M_1^3v_0/d_*^4,
 \end{split}
 \label{eq:rg-thermo-native-log-moduli}
\end{equation}
and every entrywise row-Fisher derivative difference through degree two by
\begin{equation}
 \mathcal W_B(v)=\min\{2E_B^F,
 200(H_B^F)^4[v_0+2\mathcal T_1(v)+2\mathcal T_2(v)+2\mathcal T_3(v)]\};
 \label{eq:rg-thermo-native-row-modulus-map}
\end{equation}
for two past patterns in the same $\mathfrak D$ agreeing through $Q_R$, with
$a_R=\lfloor R/2\rfloor-1$ and $\mathrm c'_n=2^n\bar k_{n+1}\mathrm c_1(1+4n)$, the
parity-jet differences obey
\begin{equation}
 \begin{split}
 V_0(B,R)&=\min\{2\rho_*^B,2\mathrm c_1B^2(R+1)e^{-\tau(R+1)}\},\\
 V_n(B,R)&=\min\bigl\{2M_n(B),B^{2n+2}\bigl[
 \mathrm c'_n(2a_R+1)^{2n}(R-a_R)e^{-\tau(R-a_R)}\\
 &\quad+2S_n(1)(a_R+2)^{2n+1}e^{-\tau(a_R+1)/n}\bigr]\bigr\},
 \qquad 1\le n\le3,
 \end{split}
 \label{eq:rg-thermo-native-past-jet-modulus}
\end{equation}
with $V_n=2M_n$ for $R=0,1$, and consequently
\begin{equation}
 \|\mathsf I^{\mathfrak D}_S-\mathsf I^{\mathfrak D}_{S'}\|_{C^2(\mathcal K),\infty,\mathrm{entry}}
 \le\omega_B(R):=\mathcal W_B(V_0,V_1,V_2,V_3),
 \label{eq:rg-thermo-native-past-fisher-modulus}
\end{equation}
which at fixed $B$ is bounded by a constant times $(R+1)^7e^{-\tau R/6}$.
Finally, the moving-expectation bound
\cref{eq:rg-spatial-proof-varying-expectation-bound} and the shell bounds
\cref{eq:rg-spatial-proof-fisher-expectation-tails} hold for the finite free
law and its finite-past rows, including targets at the free boundary, with
the modulus \cref{eq:rg-thermo-native-past-fisher-modulus}.
\status{ESTABLISHED}

\noindent\emph{Proof.}
Each assertion is the corresponding estimate of
\cref{sec:rg-spatial-proof-constrained,sec:rg-spatial-proof-higher-response,sec:rg-spatial-proof-local-jets,sec:rg-spatial-proof-moving-fisher}
for all-touching laws, and we check that each step of those proofs uses the
boundary of the finite law only through two facts. First, every sum over
incident plaquettes or over sources is a sum of nonnegative terms over a
subset of the all-touching list, so omitting the outer plaquettes only
removes terms: $\sum_in(i,e)\le6$ replaces the column identity, the path
fields keep $|a_e|\le2$, and the source lists shrink. Second, when two
conditionals of the same finite law are compared, any boundary inherited
from the original inventory, including its omitted interactions, is
identical on both sides and contributes no discrepancy.

A plaquette meets four distinct blocks and at most one edge of any retained
path, so conditional on the other fine blocks a path row is a positive
product sign law, possibly restricted to a parity fiber; the coupling of
opposite parities at Hamming distance one in
\cref{sec:rg-spatial-proof-constrained} is a property of product laws. The
tilt of a sign with baseline field $\bar h$ has total variation
$(\tanh(\bar h+j)-\tanh(\bar h-j))/2$, which is
\cref{eq:rg-thermo-native-fiber-tilt} and the bound of
\cref{eq:rg-spatial-proof-influence}, so the row costs and their telescoping
over changed plaquettes are as there. The column identity enters
\cref{eq:rg-spatial-proof-block-hamming,eq:rg-spatial-proof-weighted-column}
only as an upper bound, where $\sum_in(i,e)\le6$ suffices: each fine edge
lies in at most two plaquettes, each meeting three other blocks. Positive
finite heat baths are irreducible on the product support, and a stationary
pair coupling exists by Ces\`aro averaging on the finite simplex. Incident
anchors differ by at most one, which gives
\cref{eq:rg-thermo-native-weighted-comparison}; a shell at anchor distance
$R+1$ has at most $8|A|(R+1)$ cells and $2B^2$ fine edges per cell, which
gives the joint-target bound. For group decoupling, condition on the free
blocks outside the radius-$R$ neighborhood of $A$; by the second fact only the
newly conditioned shell contributes, and integrating over the exterior gives
\cref{eq:rg-thermo-native-group-decoupling}, which passes to infinite
constrained laws through finite-cylinder limits.

The cumulant formula and its constants are finite identities for bounded
observables, obtained by expanding the finite logarithm of a
moment-generating function, and a cumulant of two independent nonempty
groups vanishes. The source-tail bounds follow, as in
\cref{eq:rg-spatial-proof-response-tails}, from the minimum spanning tree of
the labeled anchors and a root for $A$, from
\cref{eq:rg-thermo-native-group-decoupling}, and from counting at most
$(8a_0)^rB^{2r}(M+1)^{2r}$ ordered tuples within anchor distance $M$; by the
first fact the free source list is a subset of the counted one, and
exterior-completed constant sources have zero higher cumulant. The moving
response formula is finite differentiation of a normalized finite
exponential family and holds for any finite source list.

Finite conditioning on past parities restricts the observed path fibers and
keeps every plaquette of the free law. A free target path conditioned on the
other fine blocks has independent signs with fields $h+ja_e$, $|a_e|\le2$ by
the first fact, so its parity mean has magnitude at most $\rho_*^B$ and
support in at most nine anchors, and revealing the two target parities gives
\cref{eq:rg-thermo-native-row-positivity}. The parity jets follow from
\cref{eq:rg-thermo-native-moving-response,eq:rg-thermo-native-source-tails}
as in \cref{eq:rg-spatial-proof-mean-jets}. The log amplitudes, the
row-Fisher jets and their bounds, the log-jet difference map, and the row
modulus map are algebraic identities and bounds for positive four-state rows
with this floor, as in
\cref{eq:rg-spatial-proof-log-jet-amplitudes,eq:rg-spatial-proof-log-jet-moduli,eq:rg-spatial-proof-fisher-observable-jets,eq:rg-spatial-proof-fisher-observable-bound,eq:rg-spatial-proof-fisher-observable-modulus}.
For the past modulus, split source tuples at $a_R$: near tuples occupy at
most $1+4n$ anchors in $Q_{a_R+1}$, at most $[2B^2(2a_R+1)^2]^n$ of them, and
group comparison on the remaining buffer gives the first term, the omitted
interactions producing no near discrepancy by the second fact; each far sum
obeys \cref{eq:rg-thermo-native-source-tails}. This is the argument of
\cref{eq:rg-spatial-proof-jet-modulus}, including irregular intersections with
the original inventory. The shell bounds of
\cref{sec:rg-spatial-proof-moving-fisher} use only the moving response
formula, the source tails, and the past modulus, and the shell series of a
finite past terminates. No step uses any other property of the free
boundary. $\square$

\lemmaheading{The finite-free local modulus}
{lem:rg-thermo-native-local-response}
Let a bounded parameter-independent observable $f$ have anchor support $A$
with $a_0=|A|\ge1$. Its finite-free expectations are $C^3$ on
$\widehat{\mathcal K}$. If the finite inventory contains
$A_R=\{x:\operatorname{dist}_\infty(x,A)\le R\}$ and the finite and infinite
constrained patterns agree there, then, with $\mu_{\mathsf g}$ the matching
constrained infinite law,
\begin{equation}
 \max_{|\alpha|=r}\sup_{\mathsf g\in\widehat{\mathcal K}}
 \left|\partial^\alpha\E_{\nu^{\mathfrak D}_{\mathsf g}}f-
       \partial^\alpha\E_{\mu_{\mathsf g}}f\right|
 \le\Delta_r(B,a_0,R)\|f\|_\infty,\qquad 0\le r\le3,
 \label{eq:rg-thermo-native-local-response-limit}
\end{equation}
where, for $R\ge4$,
\begin{equation}
 \Delta_0(B,a_0,R)=\min\{2,2\mathrm c_1a_0B^2(R+1)e^{-\tau(R-1)}\},
 \label{eq:rg-thermo-native-value-modulus}
\end{equation}
and, with $n_R=\lfloor R/2\rfloor-2$,
\begin{equation}
 \begin{split}
 \widetilde\Delta_r(B,a_0,R)={}&
 \bar k_{r+1}[2B^2a_0(2n_R+1)^2]^r
 \mathrm c_1(a_0+4r)B^2(R-n_R+1)e^{-\tau(R-n_R-2)}\\
 &+2U_r(B,a_0)(n_R+2)^{2r+1}e^{-\tau(n_R+1)/r},\\
 \Delta_r(B,a_0,R)={}&
 \min\{2U_r(B,a_0),\widetilde\Delta_r(B,a_0,R)\},
 \qquad 1\le r\le3,
 \end{split}
 \label{eq:rg-thermo-native-derivative-moduli}
\end{equation}
while $\Delta_r=2U_r$ for $R<4$, including $r=0$. Each modulus tends to zero
as $R\to\infty$ at fixed $B$ and $a_0$. \status{ESTABLISHED}

\noindent\emph{Proof.}
For values, condition both laws on a whole-block neighborhood one interaction
layer inside $A_R$. Every plaquette touching that smaller region is present
in both conditionals because the inventory contains $A_R$; this is the step
that uses more than \cref{lem:rg-thermo-native-free-contraction}. The
exterior assignments may have different parameter-dependent laws, but at each
$\mathsf g$ the comparison \cref{eq:rg-thermo-native-group-decoupling} is
uniform over both, and integrating it proves
\cref{eq:rg-thermo-native-value-modulus}. For derivatives, differentiate each
finite normalized law with its own source list, which gives the ordered
cumulant sums. A near source tuple and $f$ occupy at most $a_0+4r$ anchors
inside $A_{n_R+1}$, whose neighborhood of radius $R-n_R-2$ has a full extra
interaction layer inside $A_R$, so their joint total variation is at most
$\mathrm c_1(a_0+4r)B^2(R-n_R+1)e^{-\tau(R-n_R-2)}$; there are at most
$[2B^2a_0(2n_R+1)^2]^r$ near tuples, and cumulant telescoping bounds their
differences. The two far sums are bounded separately by
\cref{eq:rg-thermo-native-source-tails}. An omitted outer source cannot occur
in the common near inventory, so it is covered by the far comparison. The
infinite derivatives are those of \cref{sec:rg-spatial-proof-higher-response}:
exhausting the all-touching specifications, near tuples have uniformly
convergent finite moment laws and far sums are uniformly small, so the finite
derivatives through degree three are uniformly Cauchy and coordinate
integration identifies continuous derivatives with agreeing mixed
derivatives, also one-sidedly at $j=0$. $\square$

\lemmaheading{Local moving observables}
{lem:rg-thermo-native-moving-response}
For a local $C^2$ observable $f_{\mathsf g}$ with anchor support $A$ whose
pointwise $C^2$ jets are bounded by $E_f$, the finite and infinite
expectations of \cref{lem:rg-thermo-native-local-response} satisfy
\begin{equation}
 \|\E_{\nu^{\mathfrak D}}f-\E_\mu f\|_{C^2(\mathcal K)}
 \le E_f[\Delta_0(B,a_0,R)+2\Delta_1(B,a_0,R)+\Delta_2(B,a_0,R)].
 \label{eq:rg-thermo-native-moving-comparison}
\end{equation}
\status{ESTABLISHED}

\noindent\emph{Proof.}
By \cref{eq:rg-thermo-native-moving-response}, each derivative of order at
most two is a sum over $\beta\le\alpha$ of derivatives of order
$|\alpha-\beta|$ of expectations of the fixed observables $\partial^\beta f$.
Apply \cref{eq:rg-thermo-native-local-response-limit} to each; at degree two
the total coefficients are $1,2,1$. $\square$

\lemmaheading{Finite and infinite predictive rows}
{lem:rg-thermo-native-row-jets}
If a target has $Q_R\subset\mathfrak D$ after translation to the origin, its
finite row-Fisher observable and the infinite row-Fisher observable
$\mathsf I_{S'}$ with the same local past satisfy
\begin{equation}
 \|\mathsf I^{\mathfrak D}_S-\mathsf I_{S'}\|_{C^2(\mathcal K),\infty,\mathrm{entry}}
 \le\Omega_B(R):=\mathcal W_B\bigl(
 \Delta_0(B,1,R),\Delta_1(B,1,R),\Delta_2(B,1,R),\Delta_3(B,1,R)\bigr)
 \longrightarrow0,
 \label{eq:rg-thermo-native-free-row-comparison}
\end{equation}
pointwise in the common past coordinates. \status{ESTABLISHED}

\noindent\emph{Proof.}
Apply \cref{eq:rg-thermo-native-local-response-limit} with $a_0=1$ separately
to each target parity; when comparing the second factor the first parity is
prescribed identically on both sides. Insert the resulting jet differences in
\cref{eq:rg-thermo-native-row-modulus-map}. The bound controls row
derivatives at fixed past coordinates and does not yet differentiate an
expectation under a moving law. $\square$

\subsection{Fourth pressure derivatives and the native information tensor}
\label{sec:rg-thermo-native-pressure}

\noindent\emph{Proof of \Cref{thm:rg-thermo-native-cfour-pressure}.}
For a sign $z$, $1+\tanh(h)z=e^{hz}/\cosh h$; applying this to the $2N$ edges
and the $(L-1)^2$ plaquettes, with normalized counting in both partitions,
gives the finite line of \cref{eq:rg-thermo-native-pressure-conversion}.
The value limit of \cref{thm:rg-pressure-boundary} applies to the raw native
interaction, whose first diameter moment is at most $3t_+/(1-t_+)$ by
\cref{eq:rg-pressure-native-moment}, so
$\|\mathfrak p_L-\mathfrak p\|_{C^0(\widehat{\mathcal K})}\le12t_+/[(1-t_+)L]$,
and the conversion gives the value limit of $\mathsf A_L$. This supplies no
derivative.

Put $f_h=z_{0,1}+z_{0,2}$ and $f_j=P_0$, with anchor supports of sizes one
and three, and $m_i(\mathsf g)=\E_{\mu_{\mathsf g}}f_i$, which are $C^3$ by
\cref{lem:rg-thermo-native-local-response}. Finite differentiation gives
\begin{equation}
 \partial_h\mathsf A_L=\frac1N\sum_{x\in\Lambda_L}
                   \E_{\nu^{\mathrm f}_{L,\mathsf g}}f_h\circ\vartheta_x,
 \qquad
 \partial_j\mathsf A_L=\frac1N\sum_{x\in\{0,\ldots,L-2\}^2}
                   \E_{\nu^{\mathrm f}_{L,\mathsf g}}f_j\circ\vartheta_x,
 \label{eq:rg-thermo-native-pressure-gradient}
\end{equation}
where $\vartheta_x$ translates the observable to $x$. For $0\le r\le3$ define
\begin{equation}
 \mathcal M_r=\max\{2U_r(1,1),U_r(1,3)\},\qquad
 \bar\Delta_r(R)=\max\{2\Delta_r(1,1,R),\Delta_r(1,3,R)\}.
 \label{eq:rg-thermo-native-gradient-bounds}
\end{equation}
For $L\ge2(R+2)$ at most $4(R+2)L$ cells lack a complete target support with
its radius-$R$ neighborhood. The bulk differences of local derivatives are
bounded by $\bar\Delta_r(R)$ and every boundary derivative by $\mathcal M_r$,
an omitted plaquette being an identically zero term, so
\begin{equation}
 \max_{i\in\{h,j\},|\alpha|=r}
 \|\partial^\alpha\partial_i\mathsf A_L-
                       \partial^\alpha m_i\|_{C^0(\mathcal K)}
 \le\bar\Delta_r(R)+8\mathcal M_r\frac{R+2}{L},\qquad 0\le r\le3.
 \label{eq:rg-thermo-native-pressure-gradient-boundary}
\end{equation}
Choosing first $R$ for the local modulus and then $L$ for the boundary
fraction proves $C^3$ convergence of the gradients. On any coordinate segment
of the enlarged rectangle,
\begin{equation}
 \mathsf A_L(\mathsf g+u e_i)-\mathsf A_L(\mathsf g)
       =\int_0^u\partial_i\mathsf A_L(\mathsf g+v e_i)dv,
 \label{eq:rg-thermo-native-pressure-ftc}
\end{equation}
and uniform value and gradient convergence pass through the integral, giving
$\partial_i\mathsf A=m_i$. The same argument for the first, second, and third
gradient derivatives proves $C^4$ regularity, convergence of every derivative
through order four, and equality of mixed derivatives. Since
$|(L-1)^2/N-1|\le2/L$ is parameter independent, the conversion transfers the
result to the raw pressure. The fourth derivative uses only three
local-response derivatives, because the first derivative is already a local
source density, and no surface log-ratio estimate is differentiated. $\square$

\noindent\emph{Proof of \Cref{prop:rg-thermo-native-ctwo-metric}.}
The natural log score is
$(\mathsf T_{h,L},\mathsf T_{j,L})-\E(\mathsf T_{h,L},\mathsf T_{j,L})$, whose
covariance is $N\nabla_{\mathsf g}^2\mathsf A_L$, as for every exponential
family in its natural coordinates \citep[Ch.~3]{AmariNagaoka2000}.
\Cref{thm:rg-thermo-native-cfour-pressure} and its conversion give
\cref{eq:rg-thermo-native-fisher-pressure}. The chain rule for the normalized
log probability multiplies the source score by $D\mathsf g^\top$, which proves the
finite pullback, and $C^2$ convergence follows from the $C^2$ limit of the
source Hessian by the product and composition rules. In contrast,
\begin{equation}
 \nabla_\theta^2(\mathsf A_L\circ\mathsf g)
 =D\mathsf g^\top(\nabla_{\mathsf g}^2\mathsf A_L)D\mathsf g
   +\sum_{i=1}^2(\partial_i\mathsf A_L)\nabla_\theta^2\mathsf g_i,
 \label{eq:rg-thermo-native-nonnatural-hessian}
\end{equation}
and the last term is absent from the score covariance. $\square$

\propositionheading{A native finite-energy information floor}
{prop:rg-thermo-native-finite-energy-floor}
On an admitted source box put $H_*=h_++2j_+$, $m_*=(1+e^{2H_*})^{-1}$, and
$v_*=\operatorname{sech}^2H_*$. For $L\ge12$,
\begin{equation}
 \frac1N I(\nu^{\mathrm f}_{L,\mathsf g})\succeq
       \frac{v_*m_*^6}{8}\operatorname{Id}_2,
 \qquad \mathsf H(\mathsf g)\succeq\frac{v_*m_*^6}{8}\operatorname{Id}_2.
 \label{eq:rg-thermo-native-compact-floor}
\end{equation}
Consequently the pullback in \cref{eq:rg-thermo-native-source-pullback} has
kernel $\ker D\mathsf g$ and rank $\operatorname{rank}D\mathsf g\le2$.
\status{ESTABLISHED}

\noindent\emph{Proof.}
Every edge conditioned on all other signs has field magnitude at most $H_*$,
so either sign has conditional probability at least $m_*$ and the edge has
conditional variance at least $v_*$. Choose horizontal edges based at
$(1+3u,1+3v)$ with $u,v\ge0$ and both coordinates at most $L-2$; there are
$\lfloor L/3\rfloor^2\ge N/16$ of them for $L\ge12$, and no two share a
plaquette, so conditional on the unselected signs they are independent. For
$F=\psi_h\mathsf T_{h,L}+\psi_j\mathsf T_{j,L}$ a selected edge $e$ has
coefficient $\psi_h+\psi_jR_e$, where $R_e$ is the sum of its two neighboring
three-sign products, and
\begin{equation}
 \operatorname{Var}(F)\ge
 \E\operatorname{Var}(F\given\text{unselected signs})
 \ge v_*\sum_{e\text{ selected}}\E(\psi_h+\psi_jR_e)^2.
 \label{eq:rg-thermo-native-separated-variance}
\end{equation}
Setting the six neighboring signs of a selected edge positive forces
$R_e=2$; setting the horizontal edges immediately above and below negative
and the four vertical neighbors positive forces $R_e=-2$. Each assignment has
probability at least $m_*^6$, also under partial conditioning, because
averaging the one-edge lower bound over unspecified signs preserves $m_*$.
Therefore
\begin{equation}
 \frac{\operatorname{Var}(F)}N
 \ge\frac{v_*m_*^6}{16}
            [(\psi_h+2\psi_j)^2+(\psi_h-2\psi_j)^2]
 =\frac{v_*m_*^6}{8}(\psi_h^2+4\psi_j^2)
 \ge\frac{v_*m_*^6}{8}|\psi|^2,
 \label{eq:rg-thermo-native-directional-floor}
\end{equation}
which is the finite bound, and the source-Hessian limit inherits it. A
positive definite middle matrix in $D\mathsf g^\top\mathsf HD\mathsf g$ gives kernel
$\ker D\mathsf g$. $\square$

Only positivity of this floor is used below.

\subsection{Exact finite-free pushes and the three-cutoff limit}
\label{sec:rg-thermo-native-free-push}

\noindent\emph{Proof of \Cref{thm:rg-thermo-native-free-push-ctwo}.}
Order the cells of $\Lambda_L$ lexicographically and use the finite
predictive rows of $P^{\mathrm f}_{B,L,\mathsf g}$. Their log scores add to the
full score; each row score has conditional mean zero given its finite past,
since differentiating its normalization gives zero, and an earlier row score
is measurable for any later past, so all cross expectations vanish. With
$S_{x,L}=\Lambda_L\cap\{y:y<x\}$,
\begin{equation}
 I(P^{\mathrm f}_{B,L,\mathsf g})=
    \sum_{x\in\Lambda_L}\E_{\nu^{\mathrm f}_{BL,\mathsf g}}
                              \mathsf I^{\mathfrak D_L}_{S_{x,L}}(\mathsf g,Y),
 \label{eq:rg-thermo-native-finite-score-rows}
\end{equation}
an identity of finite $C^2$ functions, so its first two derivatives are also
exact.

Fix an inner past cutoff $n\ge0$ and an outer radius $R\ge n+4$, and consider
a bulk target $x$ with $x+Q_R\subset\Lambda_L$, translated to the origin. Its
finite past and the complete past have the same intersection with $Q_n$.
Truncating their expected row observables to that common prefix costs at most
$\epsilon_B(n)$ in entrywise $C^2$, where $W_B$, $Z^F_B$, and $\epsilon_B$
are the weight, bound, and tail of
\cref{eq:rg-spatial-proof-varying-expectation-bound,eq:rg-spatial-proof-fisher-expectation-tails},
valid for the finite free law by \cref{lem:rg-thermo-native-free-contraction}.
The finite and infinite truncated rows are functions of the same past
coordinates, with pointwise $C^2$ difference at most $\Omega_B(R)$ by
\cref{lem:rg-thermo-native-row-jets}, with $\mathfrak D_L$ the inventory of
the fine square, so their difference under the moving finite law is at most
$W_B(n)\Omega_B(R)$. It remains to compare
expectations of the same infinite truncated row under
$\nu^{\mathrm f}_{BL,\mathsf g}$ and $\mu_{\mathsf g}$; its jets are bounded by
$E_B^F$, its support is $Q_n$, and this support has a radius-$(R-n)$
neighborhood in $\mathfrak D_L$, so \cref{lem:rg-thermo-native-moving-response}
bounds the remaining cost with $a_n=(2n+1)^2$. Together,
\begin{equation}
 \begin{split}
 \mathcal B_B(n,R)={}&\epsilon_B(n)+W_B(n)\Omega_B(R)\\
 &+E_B^F[\Delta_0(B,a_n,R-n)+2\Delta_1(B,a_n,R-n)
                                  +\Delta_2(B,a_n,R-n)].
 \end{split}
 \label{eq:rg-thermo-native-three-cutoff-bulk}
\end{equation}
At most $4RL$ cells lie outside the bulk for $L\ge2R$, every finite and
infinite expected-row jet is bounded by $Z_B^F$, also near the free outer
boundary, and stationarity identifies the infinite target expectation at each
bulk translation. Dividing \cref{eq:rg-thermo-native-finite-score-rows} by $N$
gives
\begin{equation}
 \left\|N^{-1}I(P^{\mathrm f}_{B,L,\mathsf g})-\mathcal I_B(\mathsf g)
                                      \right\|_{C^2(\mathcal K),\mathrm{entry}}
 \le\mathcal B_B(n,R)+8Z_B^F\frac RL.
 \label{eq:rg-thermo-native-three-cutoff-bound}
\end{equation}
Given a tolerance, choose $n$ for the shell tail, then $R$ for the two local
comparisons at that $n$, and finally $L$ for the boundary fraction; this
order respects the growth of the constants in $a_n$, and $B$ is fixed
throughout. For infinite-law marginals the same row identity with
$\mu_{\mathsf g}$ as outer expectation gives
$\epsilon_B(R)+8Z_B^FR/L$, which is the proof of
\cref{eq:rg-spatial-proof-actual-fisher-limit}, so both limits equal
$\mathcal I_B$. At $B=1$ the channel is the identity, and
\cref{eq:rg-thermo-native-fisher-pressure,eq:rg-thermo-native-free-push-limit}
are limits of the same finite matrices, which proves
\cref{eq:rg-thermo-native-identity-density}. $\square$

\remarkheading{Area bound and native density inputs}
{rem:rg-thermo-native-push-consequences}
For every fixed $B$, $\mathcal I_B\preceq B^2\mathsf H$. If each
$\mathsf g_\gamma$ is $C^4$ on a neighborhood of a compact $\mathcal K$ with image
in a common admitted source box, the fine and coarse branch densities of a
push along a word of total factor $B$ satisfy, in $C^2(\mathcal K)$,
\[
 \frac{I(\nu^{\mathrm f}_{BL,\mathsf g_\gamma})}{B^2N}
 \longrightarrow D\mathsf g_\gamma^\top(\mathsf H\circ\mathsf g_\gamma)D\mathsf g_\gamma,
 \qquad
 \frac{I\bigl((C_B)_\#\nu^{\mathrm f}_{BL,\mathsf g_\gamma}\bigr)}{N}
 \longrightarrow
 D\mathsf g_\gamma^\top(\mathcal I_B\circ\mathsf g_\gamma)D\mathsf g_\gamma,
\]
and the per-coarse-cell Fisher loss of the push converges to
$D\mathsf g_\gamma^\top[(B^2\mathsf H-\mathcal I_B)\circ\mathsf g_\gamma]D\mathsf g_\gamma\succeq0$.
The coarse limit is a pulled-back push density, not the Hessian of a coarse
pressure. \status{ESTABLISHED}

\noindent\emph{Proof.}
\Cref{prop:rg-spatial-fisher-contraction} with blocking factor $B$ applied to
$\mathcal I_1$, and \cref{eq:rg-thermo-native-identity-density}, give the
area bound. The fine limit is \cref{prop:rg-thermo-native-ctwo-metric} on the
square of side $BL$, the coarse limit is
\cref{eq:rg-thermo-label-push-input}, and the congruence by
$D\mathsf g_\gamma$ preserves the order of the area bound. $\square$

\subsection{Why scalar boundary control is insufficient}
\label{sec:rg-thermo-native-boundary-obstruction}

\propositionheading{A bounded boundary ratio with divergent information density}
{prop:rg-thermo-native-boundary-score-counterexample}
Uniform control of log-density ratios and agreement of bulk cylinder limits
do not imply convergence of volume-normalized Fisher tensors, even for
positive smooth finite families on one carrier at each volume. At the native
product control $j=0$, the raw pressure has
$\partial_h^2\mathfrak p(h,0)=0$ while $\mathsf H_{hh}(h,0)=2\operatorname{sech}^2h>0$,
and the functions $L^{-1}\sin(L^2\theta)$ converge uniformly to zero with
derivative $L$ at zero; hence neither a raw-pressure Hessian nor a
differentiated value error supplies the Fisher density. \status{ESTABLISHED}

\noindent\emph{Proof.}
On the square carrier with $2N$ signs let $P_L$ be the fair product law,
independent of a real parameter $\theta$, and let $Q_{L,\theta}$ change only
one sign on the boundary of the centered square to the mean
$a\sin(N\theta)$, with fixed $0<a<1/2$. All probabilities are positive and
smooth on one open interval, and for every configuration
\begin{equation}
 \left|\log\frac{Q_{L,\theta}}{P_L}\right|
 \le\max\{-\log(1-a),\log(1+a)\}.
 \label{eq:rg-thermo-native-boundary-ratio-control}
\end{equation}
The varying sign leaves every fixed centered cylinder for large $L$, so both
sequences have the same bulk cylinder limit. A sign of mean $m$ has Fisher
information $(m')^2/(1-m^2)$, and independence of the other signs gives
\begin{equation}
 I(Q_{L,\theta})=
   \frac{a^2N^2\cos^2(N\theta)}{1-a^2\sin^2(N\theta)},
 \qquad I(Q_{L,0})=a^2N^2,
 \qquad N^{-1}I(Q_{L,0})=a^2N\longrightarrow\infty,
 \label{eq:rg-thermo-native-boundary-score-growth}
\end{equation}
whereas $I(P_L)=0$. The witness violates the native local jet bounds, so it
refutes the scalar boundary inference and not the native theorems. At $j=0$
the raw partition is one for every $h$, and the variance of
$\mathsf T_{h,L}$ is $2N\operatorname{sech}^2h$. $\square$

\section{Proofs for retained labels, ties, and reciprocal-volume experiments}
\label{sec:rg-thermo-label-proof}

The lattice-free statements on finite labeled families, the labeled Fisher
decomposition, the strict-phase, tie, and reciprocal-volume limits, and their
counterexamples, are in \cref{sec:cg-retained-label-asymptotics}, with the
families and norms of \cref{def:rg-thermo-label-family}; this section supplies
the native inputs, regular charts, residual scales, and comparison results
that accompany them.

\subsection{Native inputs and effective representatives}
\label{sec:rg-thermo-label-inputs}

\noindent\emph{Proof of \Cref{prop:rg-thermo-label-native-map}.}
The identity $1+(\tanh h)z=e^{hz}/\cosh h$ on every edge and plaquette gives
\cref{eq:rg-thermo-label-raw-pressure}, and the normalized branch law is the
exponential family with statistic $(\mathsf T_{h,L},\mathsf T_{j,L})$, whose
scalar factors cancel in its own normalization. Multiplying the centered
natural score by $D\mathsf g_\gamma^\top$ gives
\cref{eq:rg-thermo-label-native-pullback} and, at each finite volume, the
push identity \cref{eq:rg-thermo-label-push-pullback}. The chain and product
rules through degree two use derivatives of $\mathsf A_L$ through degree four
and of $\mathsf g_\gamma$ through degree three, and
\cref{thm:rg-thermo-native-cfour-pressure,prop:rg-thermo-native-ctwo-metric,thm:rg-thermo-native-free-push-ctwo}
give the limits, with $(L-1)^2/N\to1$ for the raw pressure. No derivative of a
$C^0$ pressure error is taken, and the word and $B$ are fixed before the
volume limit. $\square$

\propositionheading{Effective growth in both normalization modes}
{prop:rg-thermo-label-representatives}
Suppose the positive raw label weights have $C^3$ logarithms with
\begin{equation}
 \|\log a_{\gamma,L}\|_{C^3(\mathcal K)}\leq C_{\mathcal K}N,
 \qquad N^{-1}\log a_{\gamma,L}\longrightarrow\alpha_\gamma
                  \quad\text{in }C^3(\mathcal K),
 \label{eq:rg-thermo-label-raw-growth}
\end{equation}
and the branches are those of \cref{prop:rg-thermo-label-native-map}. Then
\begin{align}
 \lambda_{\gamma,L}^{\mathrm f}
 &=\log a_{\gamma,L}+\log Z_{\gamma,L},
 &\mathsf s_\gamma^{\mathrm f}&=\alpha_\gamma+\mathfrak p_\gamma,\label{eq:rg-thermo-label-factor-growth}\\
 \lambda_{\gamma,L}^{\mathrm h}
 &=\log a_{\gamma,L},
 &\mathsf s_\gamma^{\mathrm h}&=\alpha_\gamma,\label{eq:rg-thermo-label-hierarchy-growth}
\end{align}
are effective representatives of the latent experiment in the two modes, and
each satisfies \cref{eq:rg-thermo-label-effective-growth} with bulk limit
$\mathsf s_\gamma$; a fixed positive $C^3$ hierarchy $w_\gamma$ uses
$\lambda_{\gamma,L}=\log w_\gamma$ and bulk limit zero. For the all-edge
sensor of \cref{prop:rg-pressure-sensor} with mean parameter
$\eta_\gamma(\theta)$, let
\begin{equation}
 A_\gamma^{\mathrm{sens}}=\frac{1+b_\gamma\eta_\gamma}{2},\qquad
 b'_\gamma=\frac{b_\gamma+\eta_\gamma}{1+b_\gamma\eta_\gamma},
 \label{eq:rg-thermo-label-sensor-data}
\end{equation}
assume these data are $C^4$ on a neighborhood of $\mathcal K$, and assume
both the original sources and the shifted sources
$(\operatorname{atanh}b'_\gamma,j_\gamma)$ lie in common compact admitted
boxes. Then, up to a common log mass, the posterior representatives
\begin{align}
 \lambda_{\gamma,L}^{o,\mathrm f}
 &=\log a_{\gamma,L}+2N\log A_\gamma^{\mathrm{sens}}
                 +\log Z_{\gamma,L}(b'_\gamma,t_\gamma),
 \label{eq:rg-thermo-label-posterior-factor}\\
 \lambda_{\gamma,L}^{o,\mathrm h}
 &=\log a_{\gamma,L}+2N\log A_\gamma^{\mathrm{sens}}
       +\log Z_{\gamma,L}(b'_\gamma,t_\gamma)
       -\log Z_{\gamma,L}(b_\gamma,t_\gamma),
 \label{eq:rg-thermo-label-posterior-hierarchy}
\end{align}
satisfy \cref{eq:rg-thermo-label-effective-growth} with bulk limits
$\alpha_\gamma+2\log A_\gamma^{\mathrm{sens}}+\mathfrak p'_\gamma$ and
$\alpha_\gamma+2\log A_\gamma^{\mathrm{sens}}+\mathfrak p'_\gamma-\mathfrak p_\gamma$,
where $\mathfrak p'_\gamma$ is the pressure at the shifted source, and the
selected conditional branch is the shifted native free law. For a push along
a fixed word, whose label weights are those of the fine volume $B^2N$,
$\lambda_{\gamma,BL}/N\to B^2\mathsf s_\gamma$ in $C^3$ with the same strict
winner. \status{ESTABLISHED}

\noindent\emph{Proof.}
The finite formulas \cref{eq:rg-thermo-finite-log-masses} give the two
representatives, and \cref{prop:rg-thermo-label-native-map} gives the $C^3$
convergence and the bound $O(N)$ of $\log Z_{\gamma,L}$ after composition
with the $C^4$ source maps. Adding a common $c_L(\theta)$ to every
representative leaves the weights unchanged and adds $\nabla c_L$ to every
score, so their covariance is unchanged with all its derivatives; common
normalizers must therefore be removed before the growth condition is imposed,
and the bound cannot in general be imposed on $\log W_\gamma$ at a tie: for
$a_0=1$ and $a_1=e^{N\theta}$ the second derivative of $\log(a_0+a_1)$ at zero
is $N^2/4$. For the sensor, \cref{eq:rg-pressure-sensor} gives
$Z^o_{\gamma}=(A_\gamma^{\mathrm{sens}})^{2N}Z_{\gamma,L}(b'_\gamma,t_\gamma)$
by multiplication over exactly $2N$ signs, the hierarchical representative
retains its original partition divisor, and
\cref{thm:rg-thermo-native-cfour-pressure} at both sources gives the bulk
limits with all sensor derivatives. Strict positivity makes the finite
selected event well defined; the argument uses the exact source-shift identity
and applies to no other likelihood. For a push, the label weights at coarse
volume $N$ are those at fine volume $B^2N$, so
$\lambda_{\gamma,BL}/N=B^2\lambda_{\gamma,BL}/(B^2N)\to B^2\mathsf s_\gamma$,
and the winner and gaps scale by $B^2$. $\square$

\subsection{Regular charts}
\label{sec:rg-thermo-label-strict}

\propositionheading{Inverse, connection, and curvature on regular charts}
{prop:rg-thermo-label-regular-geometry}
Let symmetric tensor fields $G_L\to G$ in $C^2$ on a neighborhood of a
compact parameter chart, with $G$ positive definite there. Then
$G_L^{-1}\to G^{-1}$ in $C^2$, the Levi-Civita coefficients converge in $C^1$,
and the Riemann, Ricci, and scalar curvatures converge in $C^0$ on the chart.
For a $C^1$ family of finite laws on an open $d$-dimensional parameter domain,
with fixed positive support and positive definite Fisher tensor at a point,
$d$ of its atom probabilities form a local chart there, and distinct nearby
parameters give distinct laws. \status{ESTABLISHED}

\noindent\emph{Proof.}
Minimizing $z^\top G(\theta)z$ over the chart times the unit sphere gives a
positive minimum $\epsilon$, and uniform convergence gives
$G_L\succeq(\epsilon/2)\operatorname{Id}$ for large $L$; then
$G_L^{-1}-G^{-1}=G_L^{-1}(G-G_L)G^{-1}$ proves uniform inverse convergence.
Differentiating $G_LH_L=\operatorname{Id}$ with $H_L=G_L^{-1}$ gives
\begin{align}
 \partial_kH_L&=-H_L(\partial_kG_L)H_L,\\
 \partial_{kl}H_L
 &=H_L(\partial_lG_L)H_L(\partial_kG_L)H_L
  +H_L(\partial_kG_L)H_L(\partial_lG_L)H_L
  -H_L(\partial_{kl}G_L)H_L,
 \label{eq:rg-thermo-label-inverse-derivatives}
\end{align}
and every factor converges uniformly. With
\begin{align}
 \Gamma^m_{kl}(G_L)
 &=\frac12(G_L)^{mn}
   (\partial_k(G_L)_{ln}+\partial_l(G_L)_{kn}-\partial_n(G_L)_{kl}),\\
 R^n{}_{mkl}(G_L)
 &=\partial_k\Gamma^n_{lm}-\partial_l\Gamma^n_{km}
    +\Gamma^n_{kp}\Gamma^p_{lm}
    -\Gamma^n_{lp}\Gamma^p_{km},
 \label{eq:rg-thermo-label-connection-curvature}
\end{align}
the first line and its first derivatives converge, so the second line and its
contractions converge continuously, using metric derivatives only through
order two. For the chart, the Fisher tensor is
$(Dp)^\top\operatorname{diag}(1/p)Dp$, so the probability Jacobian has full
column rank $d$, so some $d$ atom probabilities have a nonsingular $d\times d$
Jacobian; the
inverse function theorem \citep[Thm.~9.24]{Rudin1976} makes that map
one-to-one near the point with a $C^1$ inverse. $\square$

\remarkheading{Volume rescaling of curvature}{rem:rg-thermo-curvature-rescaling}
For a constant $c>0$, replacing $G$ by $cG$ leaves the connection, the
Riemann tensor of type $(1,3)$, and the Ricci tensor unchanged, multiplies the
fully covariant Riemann tensor by $c$, and divides the scalar curvature by
$c$; in particular $N^{-1}I$ and $I$ have the same connection, and a
vanishing sequence of metrics can have divergent scalar curvature.
\status{ESTABLISHED}

\noindent\emph{Proof.}
Substitute $cG$ in \cref{eq:rg-thermo-label-connection-curvature}: the factor
$c$ cancels in $\Gamma$, lowering an index multiplies by $c$, and the scalar
contraction uses $c^{-1}G^{-1}$. $\square$

\corollaryheading{Nonempty native regular charts}
{cor:rg-thermo-native-regular-charts}
On an admitted source box, \cref{prop:rg-thermo-label-regular-geometry}
applies to $N^{-1}I(\nu^{\mathrm f}_{L,\mathsf g})\to\mathsf H$, and to a
source pullback on a compact parameter chart on which $D\mathsf g$ has full column
rank with singular values at least $\sigma_{\mathcal K}>0$. For the pushes,
write $\eta_B>0$ and $\delta_B>0$ for the floor and width of
\cref{eq:rg-spatial-proof-fisher-product-floor,eq:rg-spatial-proof-regular-fisher-region}
on a compact positive field interval. For a prescribed finite list
$\mathcal B$ of depths and a compact parameter chart whose source image lies
in that interval with $0<j<\min_{B\in\mathcal B}\delta_B$ and on which $D\mathsf g$
has full column rank with singular values at least $\sigma_{\mathcal K}>0$, the
pulled-back limits
of \cref{eq:rg-thermo-label-push-input} are bounded below by
$(\sigma_{\mathcal K}^2/2)\min_{B\in\mathcal B}\eta_B$, and
\cref{prop:rg-thermo-label-regular-geometry} applies at every listed depth
and to the strict-phase density of \cref{cor:rg-thermo-native-label-limits}
on such charts. \status{ESTABLISHED}

\noindent\emph{Proof.}
\Cref{prop:rg-thermo-native-finite-energy-floor} bounds $\mathsf H$ below on
the box, and $D\mathsf g^\top\mathsf HD\mathsf g\succeq cD\mathsf g^\top D\mathsf g\succeq
c\sigma_{\mathcal K}^2\operatorname{Id}$ for any lower bound $c$. For the
pushes, \cref{eq:rg-spatial-proof-regular-fisher-region} gives
$\mathcal I_B\succeq(\eta_B/2)\operatorname{Id}_2$ on the stated region by
\cref{thm:rg-thermo-native-free-push-ctwo}, and the same congruence gives the
common floor; the $C^2$ convergence is
\cref{prop:rg-thermo-native-ctwo-metric,prop:rg-thermo-label-native-map}, and
the strict-phase limit is a winning branch density. $\square$

The floor is not uniform in unbounded depth, and an arbitrary positive source
image need not lie in these small-$j$ charts; a direction acting only on
losing labels can disappear in the strict-phase limit, and constant rank alone
gives no quotient.

\subsection{Residual and rare-label scales}
\label{sec:rg-thermo-label-ties}

\propositionheading{Residual gradients and rare labels before score squaring}
{prop:rg-thermo-label-residual-scales}
At every finite volume the label contribution of a family of
\cref{def:rg-thermo-label-family}, with scores
$\xi_{\gamma,L}=\nabla\lambda_{\gamma,L}$, is
\begin{equation}
 \frac1N\operatorname{Cov}_{W_L}(\xi_{\gamma,L})
 =\frac1{2N}\sum_{\gamma,\gamma'}W_{\gamma,L}W_{\gamma',L}
   (\xi_{\gamma,L}-\xi_{\gamma',L})
   (\xi_{\gamma,L}-\xi_{\gamma',L})^\top.
 \label{eq:rg-thermo-label-pairwise-covariance}
\end{equation}
Consequently, if the weighted branch-density term converges, convergence of
the full density is equivalent to convergence of the right side, whose
weights cannot be replaced by their limits before the score differences are
squared. Write $\lambda_{\gamma,L}=N\mathsf s_\gamma+\varrho_{\gamma,L}$ with
$\|\varrho_{\gamma,L}\|_{C^3(\mathcal K)}/N\to0$. In the setting of
\cref{prop:rg-thermo-label-tie-scale}, suppose every label of $S$ has positive
limiting weight, fix a direction $z$ for which $z\cdot\mathsf v_\gamma$ is the
same for all $\gamma\in S$, choose $\gamma_0\in S$, and assume for some
$0\leq\upsilon<1$ that
\begin{equation}
 N^{-\upsilon}\partial_z(\varrho_{\gamma,L}-\varrho_{\gamma_0,L})(\theta_0)
 \longrightarrow d_\gamma\quad(\gamma\in S),
 \qquad d_{\gamma_0}=0.
 \label{eq:rg-thermo-label-residual-assumption}
\end{equation}
Then the directional label information divided by $N^{2\upsilon}$ tends to
$\operatorname{Var}_w(d_\gamma)$; positive variance is below the extensive
scale if $\upsilon<1/2$, gives a finite nonzero density if $\upsilon=1/2$,
and makes the density diverge if $\upsilon>1/2$. If all slopes
$\mathsf v_\gamma$ on $S$ agree, the relative residual gradients divided by
$\sqrt N$ tend to vectors $d_\gamma$, and the branch densities converge at
$\theta_0$ to $\mathcal I_\gamma(\theta_0)$, then
\begin{equation}
 \frac1NI(P_L)(\theta_0)\longrightarrow
       \sum_{\gamma\in S}w_\gamma\mathcal I_\gamma(\theta_0)
                           +\operatorname{Cov}_w(d_\gamma).
 \label{eq:rg-thermo-label-residual-density}
\end{equation}
These are limits at a fixed parameter, with no conclusion about derivatives,
and a tied label of vanishing weight requires separate control.
\status{ESTABLISHED}

\noindent\emph{Proof.}
Expanding the right side of \cref{eq:rg-thermo-label-pairwise-covariance},
the two squared terms give $N^{-1}\sum_\gamma W_\gamma\xi_\gamma\xi_\gamma^\top$
and the two cross terms give
$-N^{-1}(\sum_\gamma W_\gamma\xi_\gamma)(\sum_\gamma W_\gamma\xi_\gamma)^\top$,
which proves the identity and, with the branch term of
\cref{eq:rg-thermo-finite-labeled-fisher}, the equivalence. On $S$ equal bulk
directional slopes cancel from directional score differences at $\theta_0$;
dividing by $N^\upsilon$ and using
\cref{eq:rg-thermo-label-residual-assumption}, the limiting sum is
$\frac12\sum_{\gamma,\gamma'\in S}w_\gamma w_{\gamma'}(d_\gamma-d_{\gamma'})^2
=\operatorname{Var}_w(d_\gamma)$. Pairs with a nonmaximizing label contribute
at most $CN^{2-2\upsilon}e^{-cN}$, since all score differences are $O(N)$.
Comparing $2\upsilon$ with one gives the thresholds, and the outer-product
version with $\upsilon=1/2$ gives \cref{eq:rg-thermo-label-residual-density}.
$\square$

The subextensive and rare-label families
\cref{eq:rg-thermo-label-subextensive-example,eq:rg-thermo-label-rare-example}
of \cref{prop:rg-thermo-label-tie-counterexamples} realize every threshold of
this proposition.

\subsection{Comparison kernels and moving centers}
\label{sec:rg-thermo-label-crossover}

\propositionheading{Two parameter-independent comparison kernels}
{prop:rg-thermo-label-comparison-kernels}
Under the hypotheses of \cref{thm:rg-thermo-label-crossover}, let
$\Pi(\mathsf w)$ be the finite-label experiment on the full catalog with
masses $\bar{\mathsf p}_\gamma(\mathsf w)$, including its structural zeros,
and let $\mathcal E_L(\mathsf w)=P_{L,\theta_0+\mathsf w/N}$. There are kernels
in both directions between $\mathcal E_L(\mathsf w)$ and $\Pi(\mathsf w)$ that
are independent of $\mathsf w$ and have uniformly vanishing total-variation
errors on every compact $\mathsf w$-set. Replacing a conditional branch by its
center law costs $O(N^{-1/2})$, and the label error has the rate of
\cref{eq:rg-thermo-label-center-ratios}. \status{ESTABLISHED}

\noindent\emph{Proof.}
Let $K_L$ extract the label and let $R_L$ sample the corresponding finite
center branch:
\begin{equation}
 K_L(\gamma'\mid\gamma,x)=\boldsymbol1_{\{\gamma'=\gamma\}},\qquad
 R_L(\gamma',x\mid\gamma)=\boldsymbol1_{\{\gamma'=\gamma\}}
                                      p_{\gamma,L,\theta_0}(x).
 \label{eq:rg-thermo-label-comparison-kernels}
\end{equation}
Both are normalized, preserve labels, and depend on $L$ but not on
$\mathsf w$; the center branch is defined for every label of positive finite
mass. With
$\varepsilon_L=\sup_{\mathsf w}\frac12\sum_\gamma
 |W_{\gamma,L}(\theta_0+\mathsf w/N)-\bar{\mathsf p}_\gamma(\mathsf w)|$ over
the chosen compact, the first comparison has error at most
$\varepsilon_L\to0$. For the reverse comparison each branch has
$I(p_{\gamma,L,\theta})\preceq CN\operatorname{Id}$ along the segments in
question, by its bounded density, and on the fixed positive support
\begin{equation}
 \left\|\frac{d}{d\sigma}
        \sqrt{p_{\gamma,L,\theta_0+\sigma\mathsf w/N}}\right\|_2^2
 =\frac1{4N^2}\mathsf w^\top
       I(p_{\gamma,L,\theta_0+\sigma\mathsf w/N})\mathsf w
 \leq\frac{C|\mathsf w|^2}{4N},
 \label{eq:rg-thermo-label-square-root-bound}
\end{equation}
for $0\le\sigma\le1$, so integrating gives square-root distance at most
$C'|\mathsf w|/\sqrt N$
between the branch at $\theta_0+\mathsf w/N$ and its center. For probability
vectors, the Cauchy-Schwarz inequality gives
\begin{equation}
 \|p-p'\|_{\mathrm{TV}}
 =\frac12\sum_x|\sqrt{p(x)}-\sqrt{p'(x)}|
                      (\sqrt{p(x)}+\sqrt{p'(x)})
 \leq\|\sqrt p-\sqrt{p'}\|_2,
 \label{eq:rg-thermo-label-tv-square-root}
\end{equation}
since $\|\sqrt p+\sqrt{p'}\|_2\leq2$. Splitting
$W_\gamma p_{\gamma,\theta}-\bar{\mathsf p}_\gamma p_{\gamma,\theta_0}$ into
$W_\gamma(p_{\gamma,\theta}-p_{\gamma,\theta_0})$ and
$(W_\gamma-\bar{\mathsf p}_\gamma)p_{\gamma,\theta_0}$ on the disjoint labeled
carrier gives
\begin{equation}
 \sup_{\mathsf w}\|\mathcal E_L(\mathsf w)-\Pi(\mathsf w)R_L\|_{\mathrm{TV}}
 \leq C_{\mathcal K}N^{-1/2}+\varepsilon_L.
 \label{eq:rg-thermo-label-experiment-error}
\end{equation}
In the terminology of \citet{LeCam1964}, the two kernels bound both
deficiencies between the experiments restricted to a compact $\mathsf w$-set
by constant multiples of $\varepsilon_L$ and of the right side of
\cref{eq:rg-thermo-label-experiment-error}, so their Le Cam distance tends to
zero. $\square$

\propositionheading{The Euclidean null-slope quotient}
{prop:rg-thermo-label-quotient}
Let
\begin{equation}
 E=\operatorname{span}\{\mathsf v_\gamma-\mathsf v_{\gamma_0}:\gamma\in S\},
 \qquad r=\dim E,
 \label{eq:rg-thermo-label-slope-space}
\end{equation}
and let $\mathbb K(\mathsf w)=\operatorname{Cov}_{\bar{\mathsf p}(\mathsf w)}(\mathsf v_\gamma)$
be the window limit of \cref{thm:rg-thermo-label-crossover}. For every
$\mathsf w$, $\ker\mathbb K(\mathsf w)=E^\perp$, and
$\bar{\mathsf p}(\mathsf w)=\bar{\mathsf p}(\mathsf w')$ exactly when
$\mathsf w-\mathsf w'\in E^\perp$. The limiting family descends to
$\R^d/E^\perp$, identified with $E$, on which its probability map is
injective and, if $r>0$, its Fisher tensor is positive definite with a
positive minimum on every compact chart. The restrictions of the finite window
metrics to $E$ converge in $C^2$ with the inverse, connection, and curvature
limits of \cref{prop:rg-thermo-label-regular-geometry}. \status{ESTABLISHED}

\noindent\emph{Proof.}
All limiting weights on $S$ are positive at every finite $\mathsf w$, and the
directional quadratic form of $\mathbb K$ is their variance of
$z\cdot\mathsf v_\gamma$, which vanishes exactly when $z$ is orthogonal to all
slope differences. The probability ratios satisfy
\begin{equation}
 \log\frac{\bar{\mathsf p}_\gamma(\mathsf w)}{\bar{\mathsf p}_{\gamma_0}(\mathsf w)}
 =\log\kappa_\gamma+\mathsf w\cdot(\mathsf v_\gamma-\mathsf v_{\gamma_0}),
 \label{eq:rg-thermo-label-quotient-ratios}
\end{equation}
so equal laws are equivalent to orthogonality of $\mathsf w-\mathsf w'$ to
every slope difference, and orthogonal projection onto $E$ gives the
quotient. Choosing $r$ independent slope differences, their log ratios recover
the coordinates on $E$, which proves injectivity and a smooth inverse on the
image. In orthonormal coordinates on $E$ the restriction of $\mathbb K$ is
positive definite with a positive minimum on each compact chart, and constant
linear restriction preserves the $C^2$ convergence of
\cref{eq:rg-thermo-label-crossover-fisher}. When $r=0$ the limiting
experiment is constant and the quotient is a point. The finite metrics need
not descend along $E^\perp$, and slower scales in the other directions are not
excluded. $\square$

\propositionheading{A nonempty native realization of the limiting experiment}
{prop:rg-thermo-label-native-witness}
Every finite list of slopes $\mathsf v_\gamma\in\R^d$ and positive constants
$\kappa_\gamma$ with $\kappa_{\gamma_0}=1$ has a realization on identical
native branches with strictly positive high-temperature plaquette coupling
for which the window weights, the Fisher tensor, and both comparison kernels
above are exact at every finite volume. In this realization the experiment
that forgets the label and observes only the configuration has zero Fisher
information, so the tie scale is information about the label.
\status{ESTABLISHED}

\noindent\emph{Proof.}
Choose common $0<b_0<1$ and $0<t_0<1/12$, the free branch for every label on
one common carrier, and
$a_{\gamma,L}(\theta)=\kappa_\gamma e^{N\mathsf v_\gamma\cdot\theta}$. The
weights are positive and smooth, their logarithms have $C^3$ norms $O(N)$ on
compact sets and bulk limits $\mathsf v_\gamma\cdot\theta$, and the identical
partitions cancel in all label ratios in both modes. At $\theta_0=0$ the center
ratios are $\kappa_\gamma$ exactly, and at $\mathsf w/N$ the normalized weights
are \cref{eq:rg-thermo-label-limiting-weights}. The branches do not depend on
$\theta$, so their Fisher term is zero, the reconstruction kernel uses their
unchanged law, and both comparison errors vanish; the configuration marginal
$\sum_\gamma W_\gamma p=p$ does not depend on $\theta$ either. With two labels
in one coordinate, $\mathsf v_0=0$, $\mathsf v_1=1$, and
$\kappa_0=\kappa_1=1$,
\begin{equation}
 I_L^{(\mathsf w)}=\mathbb K(\mathsf w)=\frac1{4\cosh^2(\mathsf w/2)}>0,
 \label{eq:rg-thermo-label-binary-native-window}
\end{equation}
whose one-dimensional curvature is zero. $\square$

\theoremheading{A transverse two-label moving center}
{thm:rg-thermo-label-moving-center}
Assume the growth and branch $C^2$ hypotheses of
\cref{thm:rg-thermo-label-crossover}. Suppose exactly two labels $0,1$
maximize at $\theta_0$, all others have a strict gap there, and for a fixed
unit vector $n$, $a:=\partial_n(\mathsf s_1-\mathsf s_0)(\theta_0)>0$ after
reversing $n$ if necessary. On a sufficiently short segment through
$\theta_0$ there is, for all large $L$, exactly one center
$\theta_{c,L}=\theta_0+t_Ln$ with
$\lambda_{1,L}(\theta_{c,L})=\lambda_{0,L}(\theta_{c,L})$, and
\begin{equation}
 |t_L|\leq\frac{4|\varrho_{1,L}(\theta_0)-\varrho_{0,L}(\theta_0)|}{Na}
 \longrightarrow0,
 \qquad \varrho_{\gamma,L}=\lambda_{\gamma,L}-N\mathsf s_\gamma.
 \label{eq:rg-thermo-label-moving-center-bound}
\end{equation}
The experiment $P_{L,\theta_{c,L}+\mathsf w/N}$ has the $C^3$ weight and
$C^2$ Fisher limits of \cref{thm:rg-thermo-label-crossover} with $S=\{0,1\}$,
$\kappa_0=\kappa_1=1$, and slopes $\mathsf v_\gamma=\nabla\mathsf s_\gamma(\theta_0)$,
and its comparison kernel samples $p_{\gamma,L,\theta_{c,L}}$, independently
of $\mathsf w$. \status{ESTABLISHED}

\noindent\emph{Proof.}
Let $f=\mathsf s_1-\mathsf s_0$ and choose $r_0>0$ so that the segment
$\theta_0+[-r_0,r_0]n$ lies in the parameter neighborhood,
$\partial_nf\geq a/2$ on it, and labels $0,1$ stay uniformly above every
other bulk score there; then $f(\theta_0\pm r_0n)$ has sign $\pm$ with
magnitude at least $ar_0/2$. For $F_L=\lambda_{1,L}-\lambda_{0,L}$, $C^1$
bulk convergence gives $\partial_nF_L\geq Na/4$ on the segment for large
$L$, and uniform convergence gives opposite signs at its endpoints, so $F_L$
has exactly one zero. Since $F_L(\theta_0)=\varrho_{1,L}(\theta_0)-\varrho_{0,L}(\theta_0)=o(N)$,
integrating the derivative bound gives
\cref{eq:rg-thermo-label-moving-center-bound}. At the center the relative
mass of the two labels is one and the others are exponentially small. For
compact $\mathsf w$-sets, apply
\cref{eq:rg-thermo-label-window-increment} at the base point $\theta_{c,L}$;
compact $C^3$ growth and $\theta_{c,L}\to\theta_0$ give the limits and bounds of
\cref{eq:rg-thermo-label-window-derivatives}, and the covariance argument with
the $O(N^{-1})$ branch bound of \cref{eq:rg-thermo-label-window-decomposition}
gives $C^2$ Fisher convergence. The kernels of
\cref{eq:rg-thermo-label-comparison-kernels} with the moving center branch
and the square-root argument give the bound of
\cref{eq:rg-thermo-label-experiment-error}. $\square$

The bound permits $|t_L|$ much larger than $N^{-1/2}$, in which case sampling
the branch at $\theta_0$ would cost up to $C\sqrt N|t_L|$ in total variation;
no such replacement is used. The single transverse argument does not solve
several relative-mass equations at once, so ties of more than two labels are
not covered.

\section{Superextensive scales and prescribed actions}
\label{sec:rg-thermo-accounting-proof}

In this section the channels are coarse channels of
\cref{def:cg-coarse-channel} that preserve the label, as in
\cref{def:rg-thermo-finite-channel}, and a fixed word maps $\widetilde N=B^2N$
fine cells to $N$ coarse cells, with $B$ and the word fixed before every
volume limit. For one finite experiment and its image, $I_f$ and $I_c$ are the
full Fisher tensors in the same parameter coordinates and $D=I_f-I_c$.

\propositionheading{Preservation of a proved superextensive scale}
{prop:rg-thermo-account-superextensive}
\label{prop:rg-thermo-account-window-coordinates}
Suppose that on a compact $\mathcal K\subset\Theta$ there is $C_{\mathcal K}$,
independent of $\widetilde N$ and $\gamma$, with
\begin{equation}
 I(p_{\gamma,\widetilde N,\theta})\preceq C_{\mathcal K}\widetilde N\operatorname{Id}
 \quad\text{for all }\theta\in\mathcal K.
 \label{eq:rg-thermo-account-extensive-branch-input}
\end{equation}
Then $0\preceq D\preceq C_{\mathcal K}\widetilde N\operatorname{Id}$. If
$a_{\widetilde N}>0$ is a parameter-independent scale with
$\widetilde N/a_{\widetilde N}\to0$ and $I_f/a_{\widetilde N}\to\mathcal H$
pointwise or uniformly on $\mathcal K$, then $I_c/a_{\widetilde N}\to\mathcal H$
in the same sense; for a fixed $p>1$, $I_f/\widetilde N^p\to\mathcal H$ implies
\begin{equation}
 \frac{I_c}{N^p}\longrightarrow B^{2p}\mathcal H,
 \qquad \frac{D}{N^p}\longrightarrow0.
 \label{eq:rg-thermo-account-power-limit}
\end{equation}
For a center $\theta_{c,L}$ independent of the local coordinate, compare the
fine and coarse window coordinates
\begin{equation}
 \mathsf w=\widetilde N(\theta-\theta_{c,L}),\qquad
 \mathsf w'=N(\theta-\theta_{c,L}),\qquad \mathsf w=B^2\mathsf w'.
 \label{eq:rg-thermo-account-window-chart}
\end{equation}
In the common coordinate $\mathsf w$ the fine and coarse tensors
$\widetilde N^{-2}I_f$ and $\widetilde N^{-2}I_c$ differ by at most
$C_{\mathcal K}/\widetilde N$ uniformly on compact $\mathsf w$-sets, and if
their common limit is $\mathcal H(\mathsf w)$, the coarse tensor in its own
coordinate satisfies
\begin{equation}
 I_c^{(\mathsf w')}(\mathsf w')=B^4 I_c^{(\mathsf w)}(B^2\mathsf w')
 \longrightarrow B^4\mathcal H(B^2\mathsf w').
 \label{eq:rg-thermo-account-window-pullback}
\end{equation}
If moreover the branch densities $I(p_{\gamma,\widetilde N})/\widetilde N$ and
$I(p_{\gamma,\widetilde N}C_\gamma)/\widetilde N$ have bounded $C^2$ norms on
a compact neighborhood of the windows and the reparameterized label weights
have bounded $C^2$ norms on compact $\mathsf w$-sets, then the within-branch
loss in the $\mathsf w$ chart,
\begin{equation}
 \frac1{\widetilde N}\sum_\gamma W_{\gamma,\widetilde N}(\theta_{c,L}+\mathsf w/\widetilde N)
 \left[\frac{I(p_{\gamma,\widetilde N})-I(p_{\gamma,\widetilde N}C_\gamma)}
            {\widetilde N}\right](\theta_{c,L}+\mathsf w/\widetilde N),
 \label{eq:rg-thermo-account-window-loss}
\end{equation}
is $O(\widetilde N^{-1})$ in $C^2$, so $C^2$ convergence of the fine window
tensor transfers to the coarse one and
\cref{eq:rg-thermo-account-window-pullback} holds in $C^2$.
\status{ESTABLISHED}

\noindent\emph{Proof.}
By \cref{eq:rg-thermo-finite-within-loss},
$D=\sum_\gamma W_\gamma[I(p_\gamma)-I(p_\gamma C_\gamma)]\succeq0$, and since
each pushed branch tensor is a score Gram matrix,
$D\preceq\sum_\gamma W_\gamma I(p_\gamma)\preceq C_{\mathcal K}\widetilde N\operatorname{Id}$.
Hence $\|I_f/a_{\widetilde N}-I_c/a_{\widetilde N}\|\leq C_{\mathcal K}\widetilde N/a_{\widetilde N}\to0$
uniformly, and $\widetilde N^p=B^{2p}N^p$ gives the power case; no relation
between scales at different volumes is implied for a general sequence. The
differential of $\theta_{c,L}+\mathsf w/\widetilde N$ is
$\widetilde N^{-1}\operatorname{Id}$, which gives the factor $\widetilde N^{-2}$
and the loss bound $C_{\mathcal K}/\widetilde N$; the differential of
$\mathsf w=B^2\mathsf w'$ is $B^2\operatorname{Id}$, which gives the factor
$B^4$ and the argument $B^2\mathsf w'$. The label law at the same $\mathsf w$
is unchanged by the channel at every finite volume. For the derivative
version, every $\mathsf w$-derivative of a composed branch density contributes
a further $\widetilde N^{-1}$ and the weight derivatives are bounded, so two
applications of the product and chain rules bound the $C^2$ norm of
\cref{eq:rg-thermo-account-window-loss} by $O(\widetilde N^{-1})$. $\square$

For native branches the branch bound and the $C^2$ inputs hold by
\cref{rem:rg-thermo-native-push-consequences}. At a tie with convergent
weights and $\widetilde N^{-1}\nabla\lambda_{\gamma,\widetilde N}\to\mathsf v_\gamma$,
the score decomposition gives $I_f/\widetilde N^2\to\operatorname{Cov}_w(\mathsf v_\gamma)$,
so the coarse tensor per squared coarse area tends to
$B^4\operatorname{Cov}_w(\mathsf v_\gamma)$; in the window of
\cref{thm:rg-thermo-label-crossover} or the moving center of
\cref{thm:rg-thermo-label-moving-center}, taken in the fine area, the coarse
limit is $B^4\mathcal H(B^2\mathsf w')$, equivalently coarse slopes
$B^2\mathsf v_\gamma$. The label marginal is preserved at every finite volume,
and the discarded state information cannot contribute at these larger scales;
equality of bulk log-mass densities alone supplies neither the weight limit
nor any information limit.

\subsection{Prescribed geometry actions and relative label weights}
\label{sec:rg-thermo-account-actions}

\definitionheading{Reweighted label ensemble}{def:rg-thermo-account-reweighted}
Fix real functions $\mathscr A_{\gamma,L}(\theta)$ of the label alone and
replace the positive raw label weights by
$\widetilde a_{\gamma,L}=a_{\gamma,L}e^{-\mathscr A_{\gamma,L}}$; for a
prescribed normalized hierarchy, reweight its label masses by the same factor
and renormalize. Conditional branch laws and their physical partitions are
unchanged. A factor depending on the branch state changes the branch and is
not of this kind.

\propositionheading{Action growth and the exact relative multiplier}
{prop:rg-thermo-account-action-growth}
Suppose $\log a_{\gamma,L}$ is $C^3$ on a neighborhood of each compact
$\mathcal K\subset\Theta$ with
\begin{equation}
 \|\log a_{\gamma,L}\|_{C^3(\mathcal K)}\leq C_{\mathcal K}N,
 \qquad N^{-1}\log a_{\gamma,L}\longrightarrow\alpha_\gamma
                       \quad\text{in }C^3(\mathcal K),
 \label{eq:rg-thermo-account-raw-growth}
\end{equation}
and the actions of \cref{def:rg-thermo-account-reweighted} satisfy
\begin{equation}
 \|\mathscr A_{\gamma,L}\|_{C^3(\mathcal K)}\leq C'_{\mathcal K}N,
 \qquad N^{-1}\mathscr A_{\gamma,L}
       \longrightarrow\mathscr A_\gamma^{\mathrm{bulk}}
                       \quad\text{in }C^3(\mathcal K).
 \label{eq:rg-thermo-account-action-admission}
\end{equation}
Then the reweighted data have the same growth class, with bulk log weight
$\alpha_\gamma-\mathscr A_\gamma^{\mathrm{bulk}}$, and in either normalization
mode
\begin{equation}
 \frac{\widetilde W_{\gamma,L}}{\widetilde W_{\gamma_0,L}}
 =\frac{W_{\gamma,L}}{W_{\gamma_0,L}}
   \exp\bigl[-\mathscr A_{\gamma,L}
                   +\mathscr A_{\gamma_0,L}\bigr].
 \label{eq:rg-thermo-account-relative-action}
\end{equation}
For two identical native branches with $0<b<1$, $0<t<1/12$, equal initial
weights, and actions $\mathscr A_{0,L}=0$, $\mathscr A_{1,L}=N^\upsilon\theta$
with $0<\upsilon<1$, the categorical Fisher information at $\theta=0$ is
$N^{2\upsilon}/4$, so a vanishing bulk action can leave the extensive density
unchanged, add a finite correction, or dominate it as $\upsilon$ is below,
equal to, or above $1/2$. \status{ESTABLISHED}

\noindent\emph{Proof.}
The new log raw weight $\log a_{\gamma,L}-\mathscr A_{\gamma,L}$ keeps the
bound of order $N$ and has the stated bulk limit. The unnormalized label mass
is $a_{\gamma,L}Z_{\gamma,L}$ in the global-factor ensemble and $a_{\gamma,L}$
in a hierarchy generated from raw weights; $Z_{\gamma,L}$ is unchanged, so
either representative shifts by $-\mathscr A_{\gamma,L}$, and the ratio to
$\gamma_0$ cancels the new common normalization, also for a reweighted
prescribed hierarchy. In the example the common native law does not depend on
$\theta$, the weight ratio is $e^{-N^\upsilon\theta}$, and
\cref{eq:rg-thermo-label-binary-formula} with $f_L=-N^\upsilon\theta$ gives
$q(0)(f_L')^2=N^{2\upsilon}/4$. $\square$

An extensive action $N\mathscr A_\gamma(\theta)$ changes the bulk log-mass
density by $-\mathscr A_\gamma$, and a fixed order-one action has zero bulk
rate. Real actions define positive finite weights even where they are
negative; nonnegativity belongs to particular energy constructions, such as
\cref{rem:rg-geometry-gauge-action-sign}.

\propositionheading{Conditional limits and a bounded-action obstruction}
{prop:rg-thermo-account-action-ratios}
Let $\theta_{c,L}$ be centers at which the original positive relative weights
converge to $k_\gamma\in(0,\infty)$. If
\begin{equation}
 \mathscr A_{\gamma,L}(\theta_{c,L})-
 \mathscr A_{\gamma_0,L}(\theta_{c,L})\longrightarrow a_\gamma,
 \label{eq:rg-thermo-account-action-center-input}
\end{equation}
then the new relative weights converge to $k_\gamma e^{-a_\gamma}$, and when
these are all labels of the catalog the normalized weights converge to their
normalization. Bounded actions, even with the bulk convergence of
\cref{eq:rg-thermo-account-action-admission}, do not imply
\cref{eq:rg-thermo-account-action-center-input} or a limiting ratio.
\status{ESTABLISHED}

\noindent\emph{Proof.}
Evaluate \cref{eq:rg-thermo-account-relative-action} at the centers and use
continuity of the exponential and of normalization on a finite positive
vector; for fixed $C^3$ actions and centers tending to $\theta_0$,
$a_\gamma=\mathscr A_\gamma(\theta_0)-\mathscr A_{\gamma_0}(\theta_0)$, and a
common action cancels at every volume. For the obstruction take two identical
native branches with $0<b<1$, $0<t<1/12$, unit initial weights, and
\begin{equation}
 \mathscr A_{0,L}=0,\qquad
 \mathscr A_{1,L}=1+(-1)^L.
 \label{eq:rg-thermo-account-alternating-action}
\end{equation}
These actions are bounded by two, independent of $\theta$, and satisfy
\cref{eq:rg-thermo-account-action-admission} with zero bulk actions, yet the
ratio of label one to label zero is $e^{-2}$ for even $L$ and one for odd $L$,
so the weight of label one alternates between $e^{-2}/(1+e^{-2})$ and $1/2$.
The failure is the missing relative center limit, not a failure of
normalization or of the native branch domain. $\square$
